\documentclass[11pt,a4paper]{article}

\usepackage[utf8]{inputenc}
\usepackage[T1]{fontenc}
\usepackage[french]{babel}
\usepackage[a4paper,margin=2.2cm]{geometry}
\usepackage{amsmath,amssymb,amsthm,mathtools}
\usepackage{lmodern}
\usepackage{enumitem}
\usepackage{hyperref}
\usepackage{xcolor}
\usepackage{fancyhdr}
\usepackage{stmaryrd}
\usepackage{mathrsfs}
\usepackage{array}

\hypersetup{
    colorlinks=true,
    linkcolor=black,
    urlcolor=blue,
    citecolor=black
}

\setlength{\parindent}{0pt}
\setlength{\parskip}{0.6em}

\pagestyle{fancy}
\fancyhf{}
\fancyhead[L]{Algèbre linéaire --- Chapitre 8}
\fancyhead[R]{Applications linéaires et représentations matricielles}
\fancyfoot[C]{\thepage}

% Environnements
\newtheorem{definition}{Définition}[section]
\newtheorem{theoreme}[definition]{Théorème}
\newtheorem{proposition}[definition]{Proposition}
\newtheorem{corollaire}[definition]{Corollaire}
\newtheorem{lemme}[definition]{Lemme}
\theoremstyle{remark}
\newtheorem{remarque}[definition]{Remarque}
\newtheorem{exemple}[definition]{Exemple}
\newtheorem*{methode}{Méthode}
\newtheorem*{idee}{Idée}

% Commandes
\newcommand{\R}{\mathbb{R}}
\newcommand{\C}{\mathbb{C}}
\newcommand{\K}{\mathbb{K}}
\newcommand{\N}{\mathbb{N}}
\newcommand{\Z}{\mathbb{Z}}
\newcommand{\Vect}{\mathrm{Vect}}
\newcommand{\Ker}{\mathrm{Ker}}
\newcommand{\Img}{\mathrm{Im}}
\newcommand{\rg}{\mathrm{rg}}
\newcommand{\id}{\mathrm{Id}}
\newcommand{\tr}{\mathrm{tr}}
\newcommand{\Lin}{\mathcal{L}}
\newcommand{\Mnp}{M_{n,p}(\K)}
\newcommand{\Mpn}{M_{p,n}(\K)}
\newcommand{\Mnn}{M_n(\K)}
\newcommand{\GLn}{GL_n(\K)}

\begin{document}

\begin{center}
    {\LARGE \textbf{Chapitre 8 --- Applications linéaires et représentations matricielles}}\\[0.4em]
    {\large Version enrichie et détaillée}
\end{center}

\hrule
\vspace{1em}

\tableofcontents

\newpage

\section{Pourquoi ce chapitre est-il central ?}

Dans les chapitres précédents, on a étudié :
\begin{itemize}
    \item les espaces vectoriels ;
    \item les familles libres, génératrices et les bases ;
    \item la dimension ;
    \item les systèmes linéaires ;
    \item les matrices.
\end{itemize}

Il manque encore l’objet qui relie naturellement tous ces chapitres : les \textbf{applications linéaires}.

\begin{idee}
Une application linéaire est une transformation entre espaces vectoriels qui respecte les combinaisons linéaires.
\end{idee}

C’est pourquoi elles sont partout :
\begin{itemize}
    \item en géométrie : projections, symétries, rotations, homothéties ;
    \item en calcul : dérivation, intégration sur certains espaces ;
    \item en algèbre : passage aux matrices, changement de base, réduction ;
    \item dans les systèmes : toute matrice définit une application linéaire.
\end{itemize}

\begin{remarque}
Le vrai sens du calcul matriciel apparaît ici : une matrice n’est pas seulement un tableau de nombres, c’est souvent la représentation d’une application linéaire dans des bases choisies.
\end{remarque}

\section{Définition d’une application linéaire}

\subsection{Définition}

\begin{definition}
Soient \(E\) et \(F\) deux espaces vectoriels sur un même corps \(\K\).  
Une application
\[
u:E\to F
\]
est dite \textbf{linéaire} si
\[
\forall x,y\in E,\ \forall \lambda,\mu\in\K,\qquad
u(\lambda x+\mu y)=\lambda u(x)+\mu u(y).
\]
\end{definition}

\begin{remarque}
Cette écriture unique contient à la fois les deux propriétés usuelles :
\[
u(x+y)=u(x)+u(y)
\qquad \text{et} \qquad
u(\lambda x)=\lambda u(x).
\]
\end{remarque}

\begin{proposition}
Une application \(u:E\to F\) est linéaire si et seulement si
\[
\forall x,y\in E,\qquad u(x+y)=u(x)+u(y)
\]
et
\[
\forall \lambda\in\K,\ \forall x\in E,\qquad u(\lambda x)=\lambda u(x).
\]
\end{proposition}

\begin{proof}
Si \(u\) est linéaire au sens de la définition précédente, il suffit de prendre \((\lambda,\mu)=(1,1)\) puis \((\lambda,\mu)=(\lambda,0)\).

Réciproquement, si \(u\) vérifie les deux propriétés séparées, alors
\[
u(\lambda x+\mu y)=u(\lambda x)+u(\mu y)=\lambda u(x)+\mu u(y).
\]
\end{proof}

\subsection{Premiers exemples}

\begin{exemple}
L’application
\[
u:\R^2\to\R^2,\qquad u(x,y)=(2x-y,\ x+3y)
\]
est linéaire.
\end{exemple}

\begin{proof}
Pour \((x,y),(x',y')\in\R^2\) et \(\lambda,\mu\in\R\),
\[
u(\lambda(x,y)+\mu(x',y'))
=
u(\lambda x+\mu x',\lambda y+\mu y')
\]
\[
=
(2(\lambda x+\mu x')-(\lambda y+\mu y'),\ (\lambda x+\mu x')+3(\lambda y+\mu y')).
\]
Cela vaut
\[
\lambda(2x-y,\ x+3y)+\mu(2x'-y',\ x'+3y')
=
\lambda u(x,y)+\mu u(x',y').
\]
\end{proof}

\begin{exemple}
L’application
\[
D:\R[X]\to\R[X],\qquad P\mapsto P'
\]
est linéaire.
\end{exemple}

\begin{exemple}
L’application
\[
f:\R\to\R,\qquad x\mapsto x^2
\]
n’est pas linéaire.
\end{exemple}

\begin{proof}
En effet,
\[
f(1+1)=4
\qquad \text{tandis que} \qquad
f(1)+f(1)=2.
\]
\end{proof}

\begin{exemple}
L’application
\[
u:\R^3\to\R,\qquad u(x,y,z)=x-2y+3z
\]
est linéaire.
\end{exemple}

\section{Premières propriétés générales}

\begin{proposition}
Si \(u:E\to F\) est linéaire, alors
\[
u(0_E)=0_F.
\]
\end{proposition}

\begin{proof}
On a
\[
u(0_E)=u(0_E+0_E)=u(0_E)+u(0_E).
\]
En ajoutant l’opposé de \(u(0_E)\), on obtient
\[
u(0_E)=0_F.
\]
\end{proof}

\begin{proposition}
Si \(u:E\to F\) est linéaire, alors pour tout \(x\in E\),
\[
u(-x)=-u(x).
\]
\end{proposition}

\begin{proof}
Comme
\[
x+(-x)=0_E,
\]
on a
\[
u(x)+u(-x)=u(0_E)=0_F.
\]
Donc \(u(-x)\) est l’opposé de \(u(x)\).
\end{proof}

\begin{proposition}
Si \(u:E\to F\) est linéaire, alors pour toute combinaison linéaire
\[
\lambda_1x_1+\cdots+\lambda_px_p,
\]
on a
\[
u(\lambda_1x_1+\cdots+\lambda_px_p)
=
\lambda_1u(x_1)+\cdots+\lambda_pu(x_p).
\]
\end{proposition}

\begin{proof}
Par récurrence sur le nombre de termes.
\end{proof}

\begin{remarque}
C’est cette propriété qui fait des applications linéaires les transformations naturelles des espaces vectoriels.
\end{remarque}

\section{L’espace vectoriel \(\Lin(E,F)\)}

\subsection{Définition}

\begin{definition}
On note
\[
\Lin(E,F)
\]
l’ensemble des applications linéaires de \(E\) dans \(F\).
\end{definition}

\subsection{Structure d’espace vectoriel}

\begin{proposition}
L’ensemble \(\Lin(E,F)\), muni des opérations définies par
\[
(u+v)(x)=u(x)+v(x)
\]
et
\[
(\lambda u)(x)=\lambda\,u(x),
\]
est un espace vectoriel sur \(\K\).
\end{proposition}

\begin{proof}
Il faut vérifier que \(u+v\) est linéaire lorsque \(u\) et \(v\) le sont, de même pour \(\lambda u\), puis que les axiomes d’espace vectoriel sont satisfaits. Tout découle des propriétés de \(F\), point par point.
\end{proof}

\begin{remarque}
Cet espace vectoriel joue un rôle très important. Lorsque \(E\) et \(F\) sont de dimension finie, il sera isomorphe à un espace de matrices.
\end{remarque}

\subsection{Dimension de \(\Lin(E,F)\)}

\begin{theoreme}
Si \(E\) et \(F\) sont de dimension finie avec
\[
\dim(E)=n,\qquad \dim(F)=p,
\]
alors
\[
\dim(\Lin(E,F))=pn.
\]
\end{theoreme}

\begin{proof}
On choisit une base \(\mathcal B=(e_1,\dots,e_n)\) de \(E\) et une base \(\mathcal C=(f_1,\dots,f_p)\) de \(F\).

Pour \(1\le i\le p\), \(1\le j\le n\), on définit \(u_{ij}\in\Lin(E,F)\) par
\[
u_{ij}(e_k)=
\begin{cases}
f_i & \text{si } k=j,\\
0 & \text{si } k\neq j.
\end{cases}
\]
Cette définition a un sens et détermine une unique application linéaire.

Montrons que la famille \((u_{ij})\) forme une base de \(\Lin(E,F)\).

\textbf{Génération :} soit \(u\in \Lin(E,F)\). Écrivons
\[
u(e_j)=\sum_{i=1}^p a_{ij}f_i.
\]
Alors
\[
u=\sum_{j=1}^n\sum_{i=1}^p a_{ij}u_{ij}.
\]

\textbf{Liberté :} si
\[
\sum_{j=1}^n\sum_{i=1}^p \lambda_{ij}u_{ij}=0,
\]
en évaluant sur \(e_j\), on obtient
\[
\sum_{i=1}^p \lambda_{ij}f_i=0.
\]
Comme \((f_1,\dots,f_p)\) est libre, tous les \(\lambda_{ij}\) sont nuls.

La famille contient \(pn\) éléments, donc
\[
\dim(\Lin(E,F))=pn.
\]
\end{proof}

\section{Noyau et image}

\subsection{Définitions}

\begin{definition}
Soit \(u:E\to F\) linéaire.

Le \textbf{noyau} de \(u\), noté \(\Ker(u)\), est
\[
\Ker(u)=\{x\in E\mid u(x)=0_F\}.
\]

L’\textbf{image} de \(u\), notée \(\Img(u)\), est
\[
\Img(u)=\{u(x)\mid x\in E\}.
\]
\end{definition}

\begin{idee}
Le noyau mesure ce que l’application \textbf{annule}.  
L’image mesure ce qu’elle \textbf{produit réellement}.
\end{idee}

\subsection{Le noyau et l’image sont des sous-espaces}

\begin{proposition}
Si \(u:E\to F\) est linéaire, alors \(\Ker(u)\) est un sous-espace vectoriel de \(E\).
\end{proposition}

\begin{proof}
On a
\[
u(0_E)=0_F,
\]
donc \(0_E\in\Ker(u)\).

Si \(x,y\in\Ker(u)\) et \(\lambda,\mu\in\K\), alors
\[
u(\lambda x+\mu y)=\lambda u(x)+\mu u(y)=0_F.
\]
Donc \(\lambda x+\mu y\in\Ker(u)\).
\end{proof}

\begin{proposition}
Si \(u:E\to F\) est linéaire, alors \(\Img(u)\) est un sous-espace vectoriel de \(F\).
\end{proposition}

\begin{proof}
Comme
\[
u(0_E)=0_F,
\]
on a \(0_F\in\Img(u)\).

Si \(y_1=u(x_1)\) et \(y_2=u(x_2)\) sont dans \(\Img(u)\), alors pour \(\lambda,\mu\in\K\),
\[
\lambda y_1+\mu y_2=\lambda u(x_1)+\mu u(x_2)=u(\lambda x_1+\mu x_2)\in\Img(u).
\]
\end{proof}

\subsection{Exemples}

\begin{exemple}
Soit
\[
u:\R^2\to\R,\qquad u(x,y)=x-2y.
\]
Alors
\[
\Ker(u)=\{(x,y)\in\R^2\mid x-2y=0\}
=
\{(2t,t)\mid t\in\R\}
=
\Vect((2,1)).
\]
Et
\[
\Img(u)=\R,
\]
car tout réel \(r\) s’écrit \(r=u(r,0)\).
\end{exemple}

\begin{exemple}
Soit
\[
u:\R^3\to\R^2,\qquad u(x,y,z)=(x+y,y+z).
\]
Alors
\[
\Ker(u)=\{(x,y,z)\in\R^3\mid x+y=0,\ y+z=0\}.
\]
On obtient
\[
x=-y,\qquad z=-y,
\]
donc
\[
\Ker(u)=\{(-t,t,-t)\mid t\in\R\}=\Vect((-1,1,-1)).
\]
\end{exemple}

\section{Injectivité, surjectivité, bijectivité}

\begin{definition}
Soit \(u:E\to F\) une application.
\begin{itemize}
    \item On dit que \(u\) est \textbf{injective} si
    \[
    \forall x,y\in E,\quad u(x)=u(y)\Longrightarrow x=y.
    \]
    \item On dit que \(u\) est \textbf{surjective} si
    \[
    \forall z\in F,\ \exists x\in E,\quad u(x)=z.
    \]
    \item On dit que \(u\) est \textbf{bijective} si elle est à la fois injective et surjective.
\end{itemize}
\end{definition}

\begin{remarque}
Pour une application linéaire, ces notions sont exactement les mêmes que pour les applications entre ensembles.
La linéarité permet ensuite de les caractériser par le noyau, l’image, le rang, ou encore l’image d’une base.
\end{remarque}

\subsection{Injectivité}

\begin{theoreme}
Soit \(u:E\to F\) linéaire. Alors
\[
u \text{ injective }
\Longleftrightarrow
\Ker(u)=\{0_E\}.
\]
\end{theoreme}

\begin{proof}
Déjà donnée dans les versions précédentes : si \(u(x)=u(y)\), alors
\[
u(x-y)=0.
\]
Le caractère injectif revient donc à dire que la seule solution de \(u(x)=0\) est \(x=0\).
\end{proof}

\subsection{Surjectivité}

\begin{proposition}
Une application linéaire \(u:E\to F\) est surjective si et seulement si
\[
\Img(u)=F.
\]
\end{proposition}

\begin{proof}
C’est exactement la définition de la surjectivité.
\end{proof}

\subsection{Bijectivité}

\begin{definition}
Une application linéaire bijective est appelée un \textbf{isomorphisme}.

Un endomorphisme bijectif est appelé un \textbf{automorphisme}.
\end{definition}

\begin{definition}
Une application linéaire
\[
u:E\to E
\]
est appelée un \textbf{endomorphisme} de \(E\).
\end{definition}

\begin{remarque}
Un isomorphisme identifie deux espaces vectoriels du point de vue linéaire : ils ont exactement la même structure.
\end{remarque}

\begin{proposition}
L’inverse d’un isomorphisme est linéaire.
\end{proposition}

\begin{proof}
Soit \(u:E\to F\) un isomorphisme. Si \(y_1=u(x_1)\) et \(y_2=u(x_2)\), alors
\[
u(\lambda x_1+\mu x_2)=\lambda y_1+\mu y_2.
\]
En appliquant \(u^{-1}\), on obtient
\[
u^{-1}(\lambda y_1+\mu y_2)=\lambda u^{-1}(y_1)+\mu u^{-1}(y_2).
\]
\end{proof}

\section{Applications linéaires particulières}

\subsection{Formes linéaires}

\begin{definition}
Une application linéaire
\[
f:E\to \K
\]
est appelée une \textbf{forme linéaire}.
\end{definition}

\begin{exemple}
Sur \(\R^3\),
\[
f(x,y,z)=x-2y+z
\]
est une forme linéaire.
\end{exemple}

\subsection{Hyperplans comme noyaux de formes linéaires}

\begin{proposition}
Soit \(E\) un espace vectoriel de dimension finie \(n\), et \(f:E\to\K\) une forme linéaire non nulle. Alors
\[
\rg(f)=1
\qquad \text{et} \qquad
\dim(\Ker(f))=n-1.
\]
\end{proposition}

\begin{proof}
Comme \(f\neq 0\), son image n’est pas réduite à \(\{0\}\). Or \(\Img(f)\subset \K\), et les seuls sous-espaces de \(\K\) sont \(\{0\}\) et \(\K\). Donc
\[
\Img(f)=\K,
\]
et ainsi
\[
\rg(f)=1.
\]
Le théorème du rang donne
\[
n=\dim(\Ker(f))+1.
\]
\end{proof}

\begin{corollaire}
Le noyau d’une forme linéaire non nulle est un hyperplan.
\end{corollaire}

\subsection{Projecteurs}

\begin{definition}
Un endomorphisme \(p:E\to E\) est appelé un \textbf{projecteur} si
\[
p^2=p.
\]
\end{definition}

\begin{remarque}
Un projecteur correspond à une projection sur un sous-espace, parallèlement à un autre.
\end{remarque}

\subsection{Symétries}

\begin{definition}
Un endomorphisme \(s:E\to E\) est appelé une \textbf{symétrie} si
\[
s^2=\id_E.
\]
\end{definition}

\subsection{Homothéties}

\begin{definition}
Pour \(\lambda\in\K\), l’application
\[
h_\lambda:E\to E,\qquad x\mapsto \lambda x
\]
est une application linéaire, appelée \textbf{homothétie de rapport \(\lambda\)}.
\end{definition}

\section{Composition d’applications linéaires}

\begin{proposition}
Si
\[
u:E\to F
\qquad \text{et} \qquad
v:F\to G
\]
sont linéaires, alors
\[
v\circ u:E\to G
\]
est linéaire.
\end{proposition}

\begin{proof}
Pour \(x,y\in E\) et \(\lambda,\mu\in\K\),
\[
(v\circ u)(\lambda x+\mu y)
=
v(u(\lambda x+\mu y))
=
v(\lambda u(x)+\mu u(y))
=
\lambda v(u(x))+\mu v(u(y)).
\]
\end{proof}

\begin{remarque}
La composition des applications linéaires correspondra au produit matriciel.
\end{remarque}

\section{Image d’une famille génératrice, image d’une famille libre}

\subsection{Image d’une famille génératrice}

\begin{proposition}
Si \((e_1,\dots,e_n)\) engendre \(E\), alors
\[
(u(e_1),\dots,u(e_n))
\]
engendre \(\Img(u)\).
\end{proposition}

\begin{proof}
Si \(y\in \Img(u)\), il existe \(x\in E\) tel que \(y=u(x)\). Comme \((e_1,\dots,e_n)\) engendre \(E\),
\[
x=\lambda_1e_1+\cdots+\lambda_ne_n.
\]
Donc
\[
y=u(x)=\lambda_1u(e_1)+\cdots+\lambda_nu(e_n).
\]
\end{proof}

\subsection{Image d’une famille libre par une application injective}

\begin{proposition}
Si \(u:E\to F\) est injective et si \((e_1,\dots,e_n)\) est libre, alors
\[
(u(e_1),\dots,u(e_n))
\]
est libre.
\end{proposition}

\begin{proof}
Supposons
\[
\lambda_1u(e_1)+\cdots+\lambda_nu(e_n)=0.
\]
Alors
\[
u(\lambda_1e_1+\cdots+\lambda_ne_n)=0.
\]
Comme \(u\) est injective,
\[
\Ker(u)=\{0\},
\]
donc
\[
\lambda_1e_1+\cdots+\lambda_ne_n=0.
\]
La famille \((e_1,\dots,e_n)\) étant libre, tous les \(\lambda_i\) sont nuls.
\end{proof}

\section{Détermination par l’image d’une base}

\begin{theoreme}
Soit \((e_1,\dots,e_n)\) une base de \(E\), et soit \((y_1,\dots,y_n)\) une famille de vecteurs de \(F\). Alors il existe une unique application linéaire
\[
u:E\to F
\]
telle que
\[
u(e_i)=y_i
\qquad \text{pour tout } i.
\]
\end{theoreme}

\begin{proof}
Tout \(x\in E\) s’écrit de manière unique
\[
x=\lambda_1e_1+\cdots+\lambda_ne_n.
\]
On pose alors
\[
u(x)=\lambda_1y_1+\cdots+\lambda_ny_n.
\]
La définition est correcte car les coordonnées sont uniques. La linéarité est immédiate. L’unicité découle du fait qu’une application linéaire est entièrement déterminée par ses valeurs sur une base.
\end{proof}

\begin{corollaire}
Deux applications linéaires de \(E\) dans \(F\) qui coïncident sur une base de \(E\) sont égales.
\end{corollaire}

\section{Théorème du rang}

\subsection{Définition}

\begin{definition}
Soit \(u:E\to F\) linéaire avec \(E\) de dimension finie. On appelle \textbf{rang} de \(u\) la dimension de son image :
\[
\rg(u)=\dim(\Img(u)).
\]
\end{definition}

\subsection{Théorème}

\begin{theoreme}[Théorème du rang]
Si \(u:E\to F\) est linéaire et si \(E\) est de dimension finie, alors
\[
\dim(E)=\dim(\Ker(u))+\rg(u).
\]
\end{theoreme}

\begin{proof}
Soit
\[
(e_1,\dots,e_k)
\]
une base de \(\Ker(u)\), puis complétons-la en une base de \(E\) :
\[
(e_1,\dots,e_k,e_{k+1},\dots,e_n).
\]
On montre alors que
\[
(u(e_{k+1}),\dots,u(e_n))
\]
forme une base de \(\Img(u)\).

\textbf{Génération :} tout vecteur de \(\Img(u)\) est de la forme \(u(x)\), et \(x\) se décompose dans la base de \(E\). Les vecteurs \(e_1,\dots,e_k\) étant dans le noyau, leurs images sont nulles.

\textbf{Liberté :} si
\[
\lambda_{k+1}u(e_{k+1})+\cdots+\lambda_nu(e_n)=0,
\]
alors
\[
u(\lambda_{k+1}e_{k+1}+\cdots+\lambda_ne_n)=0.
\]
Donc cette combinaison appartient au noyau, donc à \(\Vect(e_1,\dots,e_k)\). On obtient alors une relation linéaire entre les vecteurs d’une base, donc tous les coefficients sont nuls.

Ainsi
\[
\dim(\Img(u))=n-k.
\]
Donc
\[
\dim(E)=n=k+(n-k)=\dim(\Ker(u))+\dim(\Img(u)).
\]
\end{proof}

\subsection{Conséquences}

\begin{corollaire}
Si \(u:E\to F\) est linéaire avec \(E\) de dimension finie, alors
\[
u \text{ injective } \Longleftrightarrow \rg(u)=\dim(E).
\]
\end{corollaire}

\begin{proof}
L’injectivité équivaut à \(\Ker(u)=\{0\}\), donc à \(\dim(\Ker(u))=0\), puis au théorème du rang.
\end{proof}

\begin{corollaire}
Si \(u:E\to F\) est linéaire avec \(F\) de dimension finie, alors
\[
u \text{ surjective } \Longleftrightarrow \rg(u)=\dim(F).
\]
\end{corollaire}

\section{Représentation matricielle}

\subsection{Matrice d’une application linéaire}

\begin{definition}
Soient \(E\) et \(F\) de dimensions finies, avec
\[
\mathcal B=(e_1,\dots,e_n)
\]
base de \(E\) et
\[
\mathcal C=(f_1,\dots,f_p)
\]
base de \(F\).

Pour \(u\in\Lin(E,F)\), on écrit pour chaque \(j\),
\[
u(e_j)=a_{1j}f_1+\cdots+a_{pj}f_p.
\]
La matrice
\[
A=(a_{ij})\in M_{p,n}(\K)
\]
est appelée la \textbf{matrice de \(u\) dans les bases \(\mathcal B\) et \(\mathcal C\)} et notée
\[
\mathrm{Mat}_{\mathcal B,\mathcal C}(u).
\]
\end{definition}

\begin{remarque}
La \(j\)-ème colonne de la matrice est le vecteur-colonne des coordonnées de \(u(e_j)\) dans la base \(\mathcal C\).
\end{remarque}

\subsection{Exemple}

\begin{exemple}
Soit
\[
u:\R^2\to\R^3,\qquad u(x,y)=(x+y,2x-y,3y).
\]
Dans les bases canoniques :
\[
u(1,0)=(1,2,0),\qquad u(0,1)=(1,-1,3).
\]
Donc
\[
\mathrm{Mat}(u)=
\begin{pmatrix}
1 & 1\\
2 & -1\\
0 & 3
\end{pmatrix}.
\]
\end{exemple}

\section{Lien entre application et matrice}

\begin{theoreme}
Soit \(u:E\to F\), et
\[
A=\mathrm{Mat}_{\mathcal B,\mathcal C}(u).
\]
Si \(x\in E\) a pour vecteur-colonne de coordonnées \(X\) dans \(\mathcal B\), alors \(u(x)\) a pour coordonnées
\[
Y=AX
\]
dans \(\mathcal C\).
\end{theoreme}

\begin{proof}
Écrivons
\[
x=x_1e_1+\cdots+x_ne_n.
\]
Alors
\[
u(x)=x_1u(e_1)+\cdots+x_nu(e_n).
\]
Comme les colonnes de \(A\) sont précisément les coordonnées des \(u(e_j)\), le produit \(AX\) donne exactement les coordonnées de \(u(x)\) dans \(\mathcal C\).
\end{proof}

\begin{remarque}
C’est le sens conceptuel du produit matrice-colonne.
\end{remarque}

\section{Isomorphisme entre applications linéaires et matrices}

\begin{theoreme}
Soient \(E\) et \(F\) de dimensions finies, de bases \(\mathcal B\) et \(\mathcal C\), avec
\[
\dim(E)=n,\qquad \dim(F)=p.
\]
L’application
\[
\Phi:\Lin(E,F)\to M_{p,n}(\K),\qquad
u\mapsto \mathrm{Mat}_{\mathcal B,\mathcal C}(u)
\]
est un isomorphisme d’espaces vectoriels.
\end{theoreme}

\begin{proof}
\textbf{Linéarité :} immédiate, car la matrice de \(u+v\) s’obtient en additionnant colonne par colonne celles de \(u\) et \(v\), et pareil pour \(\lambda u\).

\textbf{Injectivité :} si \(\mathrm{Mat}_{\mathcal B,\mathcal C}(u)=0\), alors toutes les colonnes sont nulles, donc
\[
u(e_j)=0
\]
pour tout \(j\). Comme \(u\) est linéaire, cela impose \(u=0\).

\textbf{Surjectivité :} soit \(A=(a_{ij})\in M_{p,n}(\K)\). Pour chaque \(j\), on pose
\[
y_j=a_{1j}f_1+\cdots+a_{pj}f_p.
\]
Le théorème de détermination par l’image d’une base fournit une unique application linéaire \(u\) telle que
\[
u(e_j)=y_j.
\]
Sa matrice est alors exactement \(A\).
\end{proof}

\begin{corollaire}
Si
\[
\dim(E)=n,\qquad \dim(F)=p,
\]
alors
\[
\dim(\Lin(E,F))=pn.
\]
\end{corollaire}

\section{Composition et matrices}

\begin{theoreme}
Si
\[
u:E\to F,\qquad v:F\to G
\]
sont linéaires, alors dans des bases compatibles,
\[
\mathrm{Mat}(v\circ u)=\mathrm{Mat}(v)\,\mathrm{Mat}(u).
\]
\end{theoreme}

\begin{proof}
Les coordonnées de \(u(x)\) sont obtenues en multipliant par \(\mathrm{Mat}(u)\), puis celles de \(v(u(x))\) en multipliant par \(\mathrm{Mat}(v)\). Donc les coordonnées de \(v\circ u\) s’obtiennent par le produit des deux matrices.
\end{proof}

\section{Endomorphismes, identité, inverse}

\subsection{Endomorphismes}

\begin{definition}
Si \(u:E\to E\) est linéaire, on parle d’\textbf{endomorphisme}.
\end{definition}

\subsection{Matrice de l’identité}

\begin{proposition}
Dans toute base \(\mathcal B\) de \(E\),
\[
\mathrm{Mat}_{\mathcal B}(\id_E)=I_n.
\]
\end{proposition}

\begin{proof}
Chaque vecteur de base est envoyé sur lui-même.
\end{proof}

\subsection{Matrice de l’inverse}

\begin{proposition}
Si \(u:E\to E\) est un automorphisme, alors
\[
\mathrm{Mat}_{\mathcal B}(u^{-1})
=
\big(\mathrm{Mat}_{\mathcal B}(u)\big)^{-1}.
\]
\end{proposition}

\begin{proof}
On utilise
\[
u\circ u^{-1}=\id_E
\qquad \text{et} \qquad
u^{-1}\circ u=\id_E,
\]
puis la compatibilité entre composition et produit matriciel.
\end{proof}

\begin{corollaire}
Un endomorphisme est un automorphisme si et seulement si sa matrice dans une base est inversible.
\end{corollaire}

\section{Noyau, image et matrice}

\subsection{Noyau}

\begin{proposition}
Soit \(u:E\to F\), et soit \(A=\mathrm{Mat}_{\mathcal B,\mathcal C}(u)\).  
Alors calculer \(\Ker(u)\) revient à résoudre le système homogène
\[
AX=0.
\]
\end{proposition}

\begin{proof}
Un vecteur \(x\in E\) appartient au noyau si et seulement si \(u(x)=0\). En coordonnées, cela équivaut à
\[
AX=0.
\]
\end{proof}

\subsection{Image}

\begin{proposition}
L’image d’une application linéaire de matrice \(A\) est engendrée par les colonnes de \(A\).
\end{proposition}

\begin{proof}
Déjà vu : si \(X=(x_1,\dots,x_n)^T\), alors
\[
AX=x_1C_1+\cdots+x_nC_n,
\]
où \(C_1,\dots,C_n\) sont les colonnes de \(A\). Donc
\[
\Img(u)=\Vect(C_1,\dots,C_n).
\]
\end{proof}

\begin{corollaire}
Le rang d’une application linéaire est le rang de sa matrice dans n’importe quelles bases.
\end{corollaire}

\section{Rang et critères matriciels}

\begin{proposition}
Soit \(u:E\to F\), avec
\[
\dim(E)=n,\qquad \dim(F)=p,
\]
et soit \(A\in M_{p,n}(\K)\) une matrice de \(u\). Alors :
\begin{enumerate}[label=\arabic*)]
    \item \(u\) est injective si et seulement si
    \[
    \rg(A)=n.
    \]
    \item \(u\) est surjective si et seulement si
    \[
    \rg(A)=p.
    \]
    \item si \(n=p\), alors \(u\) est bijective si et seulement si
    \[
    \rg(A)=n
    \]
    si et seulement si \(A\) est inversible.
\end{enumerate}
\end{proposition}

\begin{proof}
C’est la traduction matricielle des résultats généraux sur le rang, l’injectivité, la surjectivité et la bijectivité.
\end{proof}

\section{Changement de base}

\subsection{Matrice de passage}

\begin{definition}
Soient \(\mathcal B=(e_1,\dots,e_n)\) et \(\mathcal B'=(e'_1,\dots,e'_n)\) deux bases de \(E\).

La matrice dont les colonnes sont les coordonnées des \(e'_j\) dans la base \(\mathcal B\) est appelée \textbf{matrice de passage de \(\mathcal B'\) vers \(\mathcal B\)}.
\end{definition}

\begin{remarque}
Cette matrice transforme les coordonnées dans \(\mathcal B'\) en coordonnées dans \(\mathcal B\).
\end{remarque}

\subsection{Formule}

\begin{theoreme}
Soit \(u:E\to E\) un endomorphisme. Soient \(\mathcal B\) et \(\mathcal B'\) deux bases de \(E\). Si \(P\) est la matrice de passage de \(\mathcal B'\) vers \(\mathcal B\), alors
\[
\mathrm{Mat}_{\mathcal B'}(u)=P^{-1}\,\mathrm{Mat}_{\mathcal B}(u)\,P.
\]
\end{theoreme}

\begin{proof}
Si \(X'\) désigne les coordonnées de \(x\) dans \(\mathcal B'\), alors ses coordonnées dans \(\mathcal B\) sont
\[
X=PX'.
\]
Si \(A=\mathrm{Mat}_{\mathcal B}(u)\), alors les coordonnées de \(u(x)\) dans \(\mathcal B\) sont
\[
AX=APX'.
\]
Pour revenir à la base \(\mathcal B'\), on multiplie par \(P^{-1}\). On obtient donc
\[
X'_{u(x)}=P^{-1}APX'.
\]
\end{proof}

\begin{remarque}
Deux matrices représentant le même endomorphisme dans deux bases différentes sont \textbf{semblables}.
\end{remarque}

\section{Exemples classiques complets}

\subsection{Dérivation sur \(\R_2[X]\)}

\begin{exemple}
Considérons
\[
D:\R_2[X]\to \R_2[X],\qquad P\mapsto P'.
\]
Dans la base \((1,X,X^2)\),
\[
D(1)=0,\qquad D(X)=1,\qquad D(X^2)=2X.
\]
Les colonnes de la matrice sont donc :
\[
[0]_{\mathcal B}=
\begin{pmatrix}
0\\
0\\
0
\end{pmatrix},
\qquad
[1]_{\mathcal B}=
\begin{pmatrix}
1\\
0\\
0
\end{pmatrix},
\qquad
[2X]_{\mathcal B}=
\begin{pmatrix}
0\\
2\\
0
\end{pmatrix}.
\]
Donc
\[
\mathrm{Mat}_{\mathcal B}(D)=
\begin{pmatrix}
0 & 1 & 0\\
0 & 0 & 2\\
0 & 0 & 0
\end{pmatrix}.
\]
\end{exemple}

\subsection{Projecteur sur un plan}

\begin{exemple}
Considérons
\[
p:\R^3\to\R^3,\qquad p(x,y,z)=(x,y,0).
\]
Alors \(p\) est linéaire, et
\[
\Ker(p)=\Vect((0,0,1)),
\qquad
\Img(p)=\Vect((1,0,0),(0,1,0)).
\]
Dans la base canonique,
\[
\mathrm{Mat}(p)=
\begin{pmatrix}
1 & 0 & 0\\
0 & 1 & 0\\
0 & 0 & 0
\end{pmatrix}.
\]
On vérifie que
\[
p^2=p.
\]
\end{exemple}

\section{Méthodes pratiques}

\begin{methode}
Pour montrer qu’une application est linéaire, il faut tester
\[
u(\lambda x+\mu y)=\lambda u(x)+\mu u(y).
\]
\end{methode}

\begin{methode}
Pour construire la matrice d’une application linéaire :
\begin{enumerate}
    \item choisir une base de départ ;
    \item calculer l’image de chaque vecteur de cette base ;
    \item écrire chaque image dans la base d’arrivée ;
    \item placer ces coordonnées en colonnes.
\end{enumerate}
\end{methode}

\begin{methode}
Pour calculer le noyau :
\begin{enumerate}
    \item passer à la matrice ;
    \item résoudre le système homogène
    \[
    AX=0.
    \]
\end{enumerate}
\end{methode}

\begin{methode}
Pour calculer l’image :
\begin{enumerate}
    \item calculer l’image d’une base ;
    \item ou prendre les colonnes de la matrice ;
    \item puis extraire une famille libre ou calculer un rang.
\end{enumerate}
\end{methode}

\section{Erreurs classiques}

\begin{itemize}
    \item Confondre application linéaire et application affine.
    \item Oublier qu’une application linéaire doit envoyer \(0\) sur \(0\).
    \item Croire que l’image d’une famille libre est toujours libre : il faut l’injectivité.
    \item Croire que l’image d’une famille génératrice engendre tout l’espace d’arrivée : elle engendre seulement l’image.
    \item Mettre les images des vecteurs de base en lignes au lieu de les mettre en colonnes.
    \item Oublier l’ordre dans
    \[
    \mathrm{Mat}(v\circ u)=\mathrm{Mat}(v)\mathrm{Mat}(u).
    \]
    \item Se tromper dans le changement de base :
    \[
    A'=P^{-1}AP.
    \]
\end{itemize}

\section{Résumé du chapitre}

\begin{itemize}
    \item Une application linéaire respecte les combinaisons linéaires.
    \item \(\Lin(E,F)\) est un espace vectoriel.
    \item Si \(E\) et \(F\) sont de dimensions finies \(n\) et \(p\), alors
    \[
    \dim(\Lin(E,F))=pn.
    \]
    \item Le noyau et l’image d’une application linéaire sont des sous-espaces vectoriels.
    \item L’injectivité équivaut à
    \[
    \Ker(u)=\{0\}.
    \]
    \item Une application linéaire est entièrement déterminée par l’image d’une base.
    \item Le théorème du rang donne
    \[
    \dim(E)=\dim(\Ker(u))+\rg(u).
    \]
    \item Une forme linéaire non nulle a pour noyau un hyperplan.
    \item Une application linéaire entre espaces de dimension finie se représente par une matrice.
    \item Les colonnes de cette matrice sont les coordonnées des images des vecteurs de base.
    \item La composition correspond au produit matriciel.
    \item Le changement de base donne la formule
    \[
    A'=P^{-1}AP.
    \]
\end{itemize}

\section*{Exercices d’application immédiate}

\begin{enumerate}[label=\textbf{Exercice \arabic*.}]
    \item Montrer que
    \[
    u:\R^2\to\R^2,\qquad u(x,y)=(x+2y,3x-y)
    \]
    est linéaire.

    \item Déterminer le noyau et l’image de
    \[
    u:\R^2\to\R,\qquad u(x,y)=x-2y.
    \]

    \item Déterminer le noyau et l’image de
    \[
    u:\R^3\to\R^2,\qquad u(x,y,z)=(x+y,y+z).
    \]

    \item Montrer qu’une application linéaire injective envoie une famille libre sur une famille libre.

    \item Montrer qu’une application linéaire envoie une famille génératrice sur une famille génératrice de son image.

    \item Déterminer la matrice de
    \[
    u:\R^2\to\R^3,\qquad u(x,y)=(x+y,2x-y,3y)
    \]
    dans les bases canoniques.

    \item Déterminer la matrice de la dérivation sur \(\R_2[X]\) dans la base \((1,X,X^2)\).

    \item Vérifier le théorème du rang sur un exemple simple.

    \item Montrer que le noyau d’une forme linéaire non nulle est un hyperplan.

    \item Soit \(u\) un endomorphisme de \(\R^2\) de matrice
    \[
    A=
    \begin{pmatrix}
    2 & 1\\
    1 & 1
    \end{pmatrix}.
    \]
    Étudier si \(u\) est un automorphisme.

    \item Déterminer la matrice de passage entre les bases
    \[
    \mathcal B=((1,0),(0,1))
    \quad \text{et} \quad
    \mathcal B'=((1,1),(1,-1)).
    \]

    \item Relier les matrices d’un même endomorphisme dans ces deux bases.
\end{enumerate}

\section{Le dual d’un espace vectoriel}

\subsection{Définition}

\begin{definition}
Le \textbf{dual} d’un espace vectoriel \(E\) sur \(\K\) est l’espace
\[
E^\ast=\Lin(E,\K),
\]
ensemble des formes linéaires sur \(E\).
\end{definition}

\begin{proposition}
Si \(E\) est de dimension finie \(n\), alors
\[
\dim(E^\ast)=n.
\]
\end{proposition}

\begin{proof}
Comme \(E^\ast=\Lin(E,\K)\), on a en dimension finie
\[
\dim(E^\ast)=\dim(\Lin(E,\K))=\dim(E)\dim(\K)=n.
\]
\end{proof}

\subsection{Base duale}

\begin{definition}
Soit \(\mathcal B=(e_1,\dots,e_n)\) une base de \(E\). La \textbf{base duale} associée est l’unique famille
\((e_1^\ast,\dots,e_n^\ast)\) de \(E^\ast\) telle que
\[
e_i^\ast(e_j)=\delta_{ij}.
\]
\end{definition}

\begin{proposition}
La famille \((e_1^\ast,\dots,e_n^\ast)\) est une base de \(E^\ast\).
Si
\[
x=\sum_{j=1}^n x_j e_j,
\]
alors
\[
x_j=e_j^\ast(x).
\]
\end{proposition}

\subsection{Hyperplans}

\begin{definition}
Un sous-espace vectoriel \(H\subset E\) est un \textbf{hyperplan} s’il admet un supplémentaire de dimension \(1\).
En dimension finie, cela équivaut à \(\dim(H)=\dim(E)-1\).
\end{definition}

\begin{theoreme}
Si \(\varphi\in E^\ast\) est une forme linéaire non nulle et si \(E\) est de dimension finie, alors
\(\Ker(\varphi)\) est un hyperplan de \(E\).
Réciproquement, tout hyperplan est le noyau d’une forme linéaire non nulle.
\end{theoreme}

\begin{proof}
Comme \(\varphi\neq 0\), son image est un sous-espace non nul de \(\K\), donc de dimension \(1\).
Le théorème du rang donne alors \(\dim(E)=\dim(\Ker(\varphi))+1\), d’où le résultat.
La réciproque se construit en choisissant un supplémentaire de dimension \(1\).
\end{proof}

\begin{exemple}
Dans \(\R^3\), l’ensemble
\[
H=\{(x,y,z)\in\R^3\mid 2x-y+3z=0\}
\]
est un hyperplan : c’est le noyau de la forme linéaire \(\varphi(x,y,z)=2x-y+3z\).
\end{exemple}

\end{document}
